\chapter{Minimum Wages and Monopsony}\label{ch:b3:minimum-wage}

Few questions in economics have been studied as often, by as many designs, as the
effect of a minimum wage on employment. The competitive model gives a clear answer:
a binding floor raises the wage of those who keep their jobs and costs others their
jobs. A model in which each employer faces an upward-sloping supply of workers
gives another: a moderate floor can raise both wages and employment. This chapter
sets out both, then follows thirty years of evidence --- case studies, state panels,
border pairs, \omterm{def:b3:minimum-wage:bunching}{bunching estimators}, large single increases --- and states what is
agreed and what is not.

\begin{openingchart}{Five published estimates of the elasticity of employment with
respect to the minimum wage, each as stated by its authors.}{Neumark and Wascher
(1995, NBER WP 5224), for the New Jersey estimates; Dube, Lester and Reich (2010);
Neumark, Salas and Wascher (2013, NBER WP 18681)}
\begin{tikzpicture}
\begin{axis}[omrecord, xbar, bar width=9pt, width=0.62\linewidth,
  height=0.36\linewidth, xmin=-0.5, xmax=1.1, ymin=0.4, ymax=5.6,
  ytick={1,...,5}, yticklabels from table={figdata/bachelor-3/26-minimum-wage/out/elasticities.csv}{label},
  y tick label style={font=\scriptsize}, xlabel={Elasticity of employment},
  extra x ticks={0}, extra x tick labels={}, extra x tick style={grid=major}]
  \addplot[fill=omDef!60, draw=omDef] table[x=elasticity, y=pos] {figdata/bachelor-3/26-minimum-wage/out/elasticities.csv};
\end{axis}
\end{tikzpicture}
\end{openingchart}

% ledger: b3mw:nw
% ledger: b3mw:dlr
% ledger: b3mw:nsw
The same New Jersey increase of 1992 gives an elasticity of 0.93 with a telephone survey of
restaurants and $-0.24$ with their payroll records. Across all U.S. states, a panel with
state and year effects gives $-0.176$ or more in magnitude; comparing only neighbouring
counties across state borders gives 0.016. Letting the data choose the comparison
states gives about $-0.3$ for teenagers. The selection is not a sample of the literature:
it is every estimate in the studies this chapter cites that its authors state as an
elasticity with respect to the minimum wage. The spread comes from data and from the
choice of comparison, which is what the chapter is about.

\begin{recall}
The Year 1 volume built the single-buyer (monopsony) model of the labour market and
showed that a floor between the monopsony and the competitive wage raises
employment. Search and matching are the Year 2 volume's; difference in differences
is chapter 3's, regression discontinuity chapter 4's, and \omterm{def:b3:credible-empirical-work:meta-analysis}{meta-analysis} chapter 6's.
\Cref{ch:b3:labour-supply-and-demand} measured the elasticity of labour demand.
\end{recall}

\begin{notation}[Symbols of this chapter]\label{not:b3:minimum-wage:symbols}
$\bar w$ is the minimum wage; $w(L)$ the supply of labour to a firm; $\varepsilon$ its elasticity;
$\mathit{MRP}$ the marginal revenue product of labour.
\end{notation}

\section{The competitive prediction}\label{sec:b3:minimum-wage:competitive}

\begin{definition}[Own-wage elasticity of employment]\label{def:b3:minimum-wage:own-wage}
\index{own-wage elasticity of employment}
The \emph{own-wage elasticity of employment} of a minimum wage is the percentage
change in the employment of affected workers divided by the percentage change in
their average wage; the elasticity \emph{with respect to the minimum wage} divides
by the percentage change in the minimum itself.
\end{definition}

\begin{model}{The competitive labour market with a wage floor}
\assumes Employers take the wage as given; the floor binds for some workers.
\predicts Employment of affected workers falls by the labour-demand elasticity
times their wage rise; total pay of the group rises if demand is inelastic.
\tested U.S. studies of teenagers and the less educated find a preponderance of
negative estimates, more one-sided for directly affected workers than for low-wage
industries; % ledger: b3mw:ns
binding federal increases in the late 2000s, which raised the average effective
minimum by 30~\%, lowered the national employment-to-population ratio by an estimated
0.7 percentage point; % ledger: b3mw:cw
other studies find slower job growth over several years rather than an immediate
fall. % ledger: b3mw:mw
\breaks When employers set wages, a floor need not reduce employment at all.
\end{model}

\section{Monopsony and search}\label{sec:b3:minimum-wage:monopsony}

\begin{definition}[Labour supply elasticity to the firm; monopsonistic competition]\label{def:b3:minimum-wage:monopsonistic}
\index{labour supply elasticity to the firm}\index{monopsonistic competition}\index{dynamic monopsony}
The \emph{labour supply elasticity to the firm} is the percentage change in a firm's
employment when it alone raises its wage by 1~\%. In \emph{monopsonistic competition}
in the labour market, many firms each face an upward-sloping supply because
workers search, differ in tastes for jobs and quit slowly (\emph{dynamic
monopsony}).
\end{definition}

\begin{proposition}[The markdown]\label{prop:b3:minimum-wage:markdown}
A firm facing labour supply $w(L)$ with elasticity $\varepsilon$ that maximises profit pays $w = \mathit{MRP} \cdot \varepsilon/(1 + \varepsilon)$:
the wage falls short of the marginal revenue product by the share $1/(1 + \varepsilon)$. A floor between
this wage and the competitive wage raises employment.
\end{proposition}

\begin{proof}
Profit is $R(L) - w(L)L$; its derivative is $\mathit{MRP} - w - w'(L)L = \mathit{MRP} - w(1 + 1/\varepsilon)$, zero at $w = \mathit{MRP}\cdot\varepsilon/(1 + \varepsilon)$,
and $(\mathit{MRP} - w)/\mathit{MRP} = 1/(1 + \varepsilon)$. With a floor $\bar w$ above this wage, the marginal cost of a worker is
$\bar w$ up to the employment that supply provides at $\bar w$; the firm hires up to the smaller of that
supply and its demand at $\bar w$. Below the competitive wage, the supply at $\bar w$ is the smaller and
it exceeds the monopsony employment, since supply rises with the wage.
\end{proof}

\begin{omfigure}
\begin{tikzpicture}[econ, x=0.5cm, y=0.36cm]
  \draw[econ axes] (0,0) -- (12,0) node[right, font=\scriptsize] {$L$};
  \draw[econ axes] (0,0) -- (0,13) node[above, font=\scriptsize] {wage};
  \draw[curve] (0,3) -- (10,8) node[right, font=\scriptsize] {supply $w(L)$};
  \draw[curve, dashed] (0,3) -- (8.2,11.2) node[above, font=\scriptsize] {marginal cost};
  \draw[curve, omDef] (0,11) -- (11,2.2) node[right, font=\scriptsize] {MRP};
  \draw[guide] (4.44,0) -- (4.44,7.44);
  \draw[guide] (0,5.22) -- (4.44,5.22);
  \draw[curve, omThm] (0,5.8) -- (5.6,5.8);
  \node[above, font=\scriptsize] at (1.4,5.8) {floor $\bar w$};
  \draw[guide] (5.6,0) -- (5.6,5.8);
  \node[below, font=\scriptsize] at (4.44,0) {$L^m$};
  \node[below, font=\scriptsize] at (5.8,0) {$L^{\bar w}$};
  \node[left, font=\scriptsize] at (0,5.22) {$w^m$};
\end{tikzpicture}
\par\smallskip{\small The single buyer hires where marginal cost meets the marginal
revenue product, $L^m$, and pays $w^m$ on the supply curve. A floor $\bar w$ makes the marginal cost
flat at $\bar w$ up to the supply curve: employment rises to $L^{\bar w}$.}
\end{omfigure}

With many firms, the search model of the Year 2 volume gives the same upward slope:
when workers find jobs one at a time and quit to better offers, a firm that pays
more fills vacancies faster and loses fewer workers. In the Burdett--Mortensen model
each firm faces an upward-sloping supply even though it is one of many; the result
is stated here without proof.

\begin{model}{Monopsony in the labour market}
\assumes Each employer faces an upward-sloping supply of workers.
\predicts Wages are marked down below the marginal product; a moderate floor raises
wages without lowering employment, and can raise it.
\tested On an online task platform, observational and experimental estimates of
the supply elasticity to an employer are around 0.1, a large markdown; % ledger: b3mw:djns
across 8{,}000 U.S. occupation-by-area labour markets, moving from the 25th to the 75th percentile
of employer concentration goes with posted wages 17~\% lower. % ledger: b3mw:ams
\breaks A floor above the competitive wage reduces employment in this model too;
how far a given floor is from that point is the empirical question.
\end{model}

\begin{definition}[Kaitz index]\label{def:b3:minimum-wage:kaitz}
\index{Kaitz index}
The \emph{Kaitz index} is the ratio of the minimum wage to the median wage.
\end{definition}

\begin{dated}{2026}
The United Kingdom sets its adult minimum by a \omterm{def:b3:minimum-wage:kaitz}{Kaitz index}: the rate is not to fall below
two-thirds of median hourly earnings. The April 2024 rate of 11.44 pounds an hour fell
short --- 11.76 would have been needed --- so its index was $11.44/(1.5 \times 11.76) = 0.65$; the Low
Pay Commission recommended 12.21 pounds for April 2025, and 1.75 million jobs were paid
at the minimum in 2024. % ledger: b3mw:lpc.bite
The United States federal minimum has been 7.25 dollars an hour since 24 July 2009,
the latest change in the Labor Department's record, so its index falls whenever the
median wage rises; the New Jersey episode below was a state raising its own floor
above the federal one. % ledger: b3mw:dol.fed
\end{dated}

\section{The first natural experiment and the replies}\label{sec:b3:minimum-wage:nj}

\begin{record}[New Jersey and Pennsylvania, 1992]
On 1 April 1992 New Jersey raised its minimum from 4.25 to 5.05 dollars an hour; Pennsylvania's
stayed at 4.25. A survey of 410 fast-food restaurants before and after found employment
in New Jersey up 13~\% relative to Pennsylvania, with high-wage New Jersey stores, which
the law did not touch, growing like those in Pennsylvania. % ledger: b3mw:ck
Payroll records from 230 restaurants of four chains gave the opposite sign: the
standard deviation of employment changes in the telephone survey was three times that in the
payroll data, and the payroll estimate was a fall of 4.6~\%. % ledger: b3mw:nw
A reply using the government's employment records for all restaurants found
similar or slightly faster growth in New Jersey, and no relative loss in Pennsylvania
when the federal minimum later rose there; the payroll sample's difference was traced
to the restaurants of a single franchisee. % ledger: b3mw:ck2
\end{record}

\section{Border pairs and bunching}\label{sec:b3:minimum-wage:designs}

\begin{definition}[Border-pair design]\label{def:b3:minimum-wage:border-pair}
\index{border-pair design}
A \emph{border-pair design} compares neighbouring counties on opposite sides of a
state border with different minimum wages, using each pair's common shocks as the
control.
\end{definition}

\begin{definition}[Bunching estimator; spillover effect]\label{def:b3:minimum-wage:bunching}
\index{bunching estimator}\index{spillover effect}
A \emph{bunching estimator} measures the employment effect of a minimum wage as the
number of excess jobs paying at or just above the new minimum less the number of
jobs missing below it. A \emph{spillover effect} is a rise in wages just above the
minimum caused by the minimum.
\end{definition}

\begin{proposition}[Missing and excess jobs]\label{prop:b3:minimum-wage:bunching}
If a minimum wage affects only jobs below some wage $\bar w + k$, the change in employment
equals the change in the number of jobs below $\bar w$ plus the change in the number between
$\bar w$ and $\bar w + k$, each measured against a control.
\end{proposition}

\begin{proof}
Employment is the sum of the jobs in every wage bin. Jobs above $\bar w + k$ are unaffected,
so their change equals the control's; the effect on total employment is the effect
on the bins below $\bar w + k$, which is the sum over the two ranges. The jobs below $\bar w$
either vanish or move up; the excess in the window counts those that moved, plus any
spillover from bins in the window.
\end{proof}

\begin{omchart}{Brenland's working heads by hourly wage (one-dalen bins), before a
floor of 12 dalens and after it under the problem's assumption that a tenth of the jobs
below 12 disappear and the rest are paid 12--13.}{Invented data: Brenland's household
survey; the weekend problem's assumption}
\begin{tikzpicture}
\begin{axis}[omrecord, ybar, bar width=3pt, width=0.66\linewidth,
  height=0.36\linewidth, xmin=2, xmax=31, ymin=0, ymax=380,
  xlabel={Hourly wage (dalens)}, ylabel={Heads}, legend pos=north east,
  legend style={font=\footnotesize}]
  \addplot[fill=omDef!50, draw=omDef, area legend] table[x=bin, y=before] {figdata/bachelor-3/26-minimum-wage/out/brenland_hist.csv};
  \addplot[fill=omThm!60, draw=omThm, area legend] table[x=bin, y=after] {figdata/bachelor-3/26-minimum-wage/out/brenland_hist.csv};
  \legend{Before, After}
\end{axis}
\end{tikzpicture}
\end{omchart}

% figdata: brenland_hist
Applied to 138 state minimum-wage increases in the United States from 1979 to 2016,
the method finds that the number of low-wage jobs was essentially unchanged over the
five years after each increase, with modest spillovers above the new minimum and
some job loss in sectors whose goods are traded across states. % ledger: b3mw:cdlz
Border pairs give the same answer for restaurants: an elasticity of 0.016, ruling out
values below $-0.147$ at the 90~\% level. % ledger: b3mw:dlr

\section{Other margins}\label{sec:b3:minimum-wage:margins}

Employment is not the only margin. Firms can cut hours, raise prices, replace
workers with machines, or close. Hungary's large and lasting increase gives
evidence on all of them: four years on, employment elasticities were negative but
small; consumers paid about 75~\% of the cost through higher prices and owners 25~\%; firms
substituted capital for labour; and job losses were larger where prices were harder
to raise. % ledger: b3mw:hl
In Seattle, an increase to 13 dollars in 2016 cut hours in low-wage jobs by 6--7~\% while
their hourly wages rose 3~\%, so that pay per job fell by 74 dollars a month; employment
in restaurants at all wage levels showed no effect. % ledger: b3mw:seattle

\begin{record}[Germany's national minimum wage, 2015]
Germany introduced a national minimum wage that affected 15~\% of employees. Wages rose
and employment did not fall. Low-wage workers moved from smaller, lower-paying and
less productive establishments to larger, better-paying and more productive ones,
and this upgrading accounts for up to 17~\% of the wage increase. % ledger: b3mw:dust
\end{record}

\section{Small and large increases}\label{sec:b3:minimum-wage:size}

\begin{debate}[The employment effect of a minimum wage]
\emph{The question.} How much employment do minimum wages cost, at the levels
actually chosen?
\emph{State panels with fixed effects} find negative effects for teenagers, near $-0.3$
with data-chosen controls (\emph{Grade 2}; weakness: the choice of control states and
trends drives the sign). % ledger: b3mw:nsw
\emph{Border pairs} find about zero for restaurants (\emph{Grade 2}; weakness: spillovers
across the border, and neighbours that differ in trends). % ledger: b3mw:dlr
\emph{\omterm{def:b3:minimum-wage:bunching}{Bunching estimators}} find the count of low-wage jobs unchanged (\emph{Grade 2};
weakness: the window above the minimum is chosen by the researcher). % ledger: b3mw:cdlz
\emph{Large single increases}: Seattle cut hours (\emph{Grade 2}; weakness: one city
and a short window); % ledger: b3mw:seattle
Hungary's increase lowered employment a little and raised prices (\emph{Grade 2}); % ledger: b3mw:hl
Germany's raised wages without lowering employment (\emph{Grade 2}). % ledger: b3mw:dust
\emph{What is agreed}: moderate increases have small effects on the employment of
the groups studied, and much of the cost is passed on in prices. \emph{What is
open}: the effect of increases that take the minimum high relative to the median
wage, where the competitive and the monopsony models both predict losses and
differ on where they start.
\end{debate}

\section{Exercises}\label{sec:b3:minimum-wage:exercises}

\begin{exercise}[$\star$]\label{exo:b3:minimum-wage:1}
In the opening chart, how many estimates are negative? What separates the two New
Jersey estimates?
\end{exercise}

\begin{exercise}[$\star$]\label{exo:b3:minimum-wage:2}
With an elasticity of employment with respect to the minimum of $-0.2$, what does a 10~\% rise
in the minimum do to employment of the affected group?
\end{exercise}

\begin{exercise}[$\star$]\label{exo:b3:minimum-wage:3}
A firm's labour supply is $w = 4 + 0.2L$ and its marginal revenue product $24 - 0.2L$. Find the
monopsony employment and wage, and the competitive ones.
\end{exercise}

\begin{exercise}[$\star$]\label{exo:b3:minimum-wage:4}
What observation would contradict monopsony power for a given employer?
\end{exercise}

\begin{exercise}[$\star\star$]\label{exo:b3:minimum-wage:5}
A firm's supply elasticity is 4. By what share of the marginal revenue product is the
wage marked down?
\end{exercise}

\begin{exercise}[$\star\star$]\label{exo:b3:minimum-wage:6}
In exercise 3, find employment with a floor of 12, of 14 and of 16.
\end{exercise}

\begin{exercise}[$\star\star$]\label{exo:b3:minimum-wage:7}
An invented histogram of jobs by wage bin is given before and after a minimum rises
from 8 to 10, in treated and control states. Compute the missing jobs, the excess jobs
from 10 to 13 and the employment change.
\end{exercise}

\begin{exercise}[$\star\star$]\label{exo:b3:minimum-wage:8}
Labour is 30~\% of a restaurant's costs; a minimum wage raises its wage bill by 10~\%. With
full pass-through, by how much do its prices rise?
\end{exercise}

\begin{exercise}[$\star\star$]\label{exo:b3:minimum-wage:9}
Chapter 6 defined the \omterm{def:b3:credible-empirical-work:meta-analysis}{funnel plot}. If the minimum-wage elasticities of the
literature were plotted against their precision, what asymmetry would suggest
selective publication? Why could the five estimates of the opening chart not be
used to draw one?
\end{exercise}

\begin{exercise}[$\star\star\star$]\label{exo:b3:minimum-wage:10}
Prove the markdown formula, and show that the markdown vanishes as the firm's supply
becomes perfectly elastic.
\end{exercise}

\begin{exercise}[$\star\star\star$]\label{exo:b3:minimum-wage:11}
Five hundred invented firms face the supply $L = 10w^4$ and differ in productivity. Estimate
the supply elasticity by regressing log employment on log wage, and the implied
markdown.
\end{exercise}

\begin{exercise}[$\star\star\star$]\label{exo:b3:minimum-wage:12}
A region raises its minimum by 40~\%, to a \omterm{def:b3:minimum-wage:kaitz}{Kaitz index} of 0.8. Which model applies, and which
evidence of the chapter is most relevant?
\end{exercise}

\section{Problem: A Minimum Wage for Brenland}\label{sec:b3:minimum-wage:problem}

\begin{problem}[{Weekend problem --- a minimum wage for Brenland}]\label{pb:b3:minimum-wage:brenland}
Brenland considers a floor of 12 dalens an hour. Use the hourly wages (gross income over
hours) of the working heads in the survey (\code{load("households")}). A typical
Brenland firm faces labour supply $w = 6 + 0.1L$ and has a marginal revenue product of
$20 - 0.1L$, with $L$ workers. % figdata: brenland_stats, brenland_hist

\medskip\noindent\textbf{Part I --- the competitive prediction.}
\begin{enumerate}
\item Find the median hourly wage and the \omterm{def:b3:minimum-wage:kaitz}{Kaitz index} of the floor.
\item What share of working heads earn less than 12?
\item By how much, on average, does the floor raise their wages?
\item With a labour-demand elasticity of 1.0, predict the employment loss among them.
\item Why is a prediction from one elasticity unreliable for wage rises this large?
\end{enumerate}

\medskip\noindent\textbf{Part II --- the firm with market power.}
\begin{enumerate}[resume]
\item Find the firm's monopsony employment and wage.
\item Find its marginal revenue product there, and the share of it the wage falls
short by.
\item Compute the supply elasticity at that point, and check the markdown formula.
\item Find the competitive employment and wage.
\item Find employment with floors of 12, 15 and 17.
\item What floor maximises employment?
\item Above what floor does employment fall below the monopsony level?
\end{enumerate}

\medskip\noindent\textbf{Part III --- bunching.} After the floor, a tenth of the jobs
below 12 disappear and the rest are paid 12--13.
\begin{enumerate}[resume]
\item Count the missing jobs in the survey.
\item Count the excess jobs at 12--13.
\item Compute the bunching estimate of the employment change, in heads and as a share
of working heads.
\item What would spillovers to the 13--14 bin change in the estimate?
\item How would you check the assumption that jobs above 13 are unaffected?
\end{enumerate}

\medskip\noindent\textbf{Part IV --- a border comparison.} Kestria has no floor. In the
regions along the border, employment goes from 100 to 99.0 on Brenland's side and from 100
to 100.5 on Kestria's.
\begin{enumerate}[resume]
\item Compute the difference-in-differences estimate.
\item By how much does the floor raise the average wage of all working heads?
\item Compute the implied \omterm{def:b3:minimum-wage:own-wage}{own-wage elasticity of employment}.
\item Compare it with the competitive prediction of Part I.
\item Name two threats to the border comparison.
\item How could workers crossing the border bias the estimate?
\item What would the result suggest about Brenland's employers?
\item How would you design a better test before choosing a higher floor?
\item State the result: the employment-maximising minimum wage in dalens an hour, and the
measured employment effect of the actual floor.
\end{enumerate}
\end{problem}
